% Source-derived chapter 09: Further results
% Original source: riemann-roch-polynomials-hyperkahler-fourfolds
\chapter{Further results}
\label{riemann-roch-polynomials-hyperkahler-fourfolds-chapter-09}
\source{riemann-roch-polynomials-hyperkahler-fourfolds}

\begin{remark}[Source section]
\label{riemann-roch-polynomials-hyperkahler-fourfolds-chapter-09-source}
\source{riemann-roch-polynomials-hyperkahler-fourfolds}
This chapter records the source section; no Lean declaration has yet been checked for this material.
\notready
\end{remark}

% The source section title was: Further results
\label{sec:further}
\source{riemann-roch-polynomials-hyperkahler-fourfolds}
We keep the same setup: $X$ is a \hk\ manifold of dimension $2n$ with classes $\lll,\mm\in H^2(X,\Z)$ such that $\int_X\lll^{2n}=0$ and
$$a\coloneqq \frac 1{n!} \int_X\lll^n\mm^n$$
is a nonzero integer (Lemma~\ref{lemmaa});  changing $\mm$ into $-\mm$ if necessary, we may assume that the integers~$q_X(\lll,\mm)$, hence also $a$, are positive.\ After dealing for most of   this article  with the case~$a=1$, we examine in this section the general case.

\subsection{Boundedness results}

We first prove a general boundedness result, assuming $n$ and $a$ are fixed (this is~\cite[Theorem~4.9]{kam}).

\begin{prop}
\source{riemann-roch-polynomials-hyperkahler-fourfolds}
\notready\label{prop41}
Let $a$ be a fixed positive integer.\
The number of deformation types of \hk\ manifolds $X$ of fixed dimension $2n$ for which there are classes~$\lll,\mm\in H^2(X,\Z)$ such that $\int_X \lll^{2n}=0$ and $\int_X \lll^{n}\mm^{n}=an!$ is finite.
\end{prop}

\begin{proof}
As explained above, we may assume $q_X(\lll,\mm)>0$.\
Since we can always add to $\mm$ a  multiple of $\lll$ (which changes $q_X(\mm)$ by adding multiples of $2q_X(\lll,\mm)$ but neither $q_X(\lll,\mm)$ nor~$a$), we may further assume
\begin{equation}\label{eqq}
\source{riemann-roch-polynomials-hyperkahler-fourfolds}
-q_X(\lll,\mm)<q_X(\mm)\le q_X(\lll,\mm).
\end{equation}

We have, using~\eqref{fuj}, \eqref{eqq}, and \eqref{defia},
\begin{equation*}
\int_X(\lll+\mm)^{2n}=c_X q_X(\lll+\mm)^n =c_X(2 q_X(\lll,\mm)+q_X(\mm) )^n\le c_X 3^n q_X(\lll,\mm)^n= a 3^n\frac{(2n)!}{2^nn!}  .
\end{equation*}
By the surjectivity of the period map, there exists a deformation of $X$ whose N\'eron--Severi group is generated by the class $\lll+\mm$, which is therefore ample (or antiample).\
One can then apply~\cite[Corollary~1.2]{huyf} to conclude.
\end{proof}



 In dimension~$4$, we improve this result by allowing $a$ to take infinitely many values.\


\begin{prop}
\source{riemann-roch-polynomials-hyperkahler-fourfolds}
\notready\label{prop42}
The number of deformation types of  \hk\ fourfolds $X$   for which there are classes~$\lll, \mm\in H^2(X,\Z)$ such that   \mbox{$\int_X \lll^4=0$} and  the   integer $a= \frac{1}{2}\int_X \lll^{2}\mm^{2}$ is positive and square-free, is finite.\end{prop}



\begin{proof}
The proposition is a consequence of Lemma~\ref{lemmax}, whose notation we keep.\ Indeed, that lemma and~\eqref{ax}
 say that the integer $aA'_X\coloneqq 12^2 \cdot 2aA_X$ is a perfect (positive) square.\
Since~$a$ is square-free and $A'_X\le 262$ by Lemma~\ref{lemmguan}, the integer $a$ is a product of (distinct) prime numbers that are all $<262$.\
It is therefore bounded (by the product of all these primes) and boundedness for $X$ follows from Proposition~\ref{prop41}.
\end{proof}


\subsection{Low values of $a$ for \hk\ fourfolds}

We obtain restrictions on the values that the integer $a$ may take for \hk\ fourfolds.\ The case $a=1$ was analyzed in Theorem~\ref{th43a} and is completely clarified by Theorem~\ref{thm:MainThm}:   it is only realized by \hk\ manifolds of~$\mathrm{K3}^{[2]}$ deformation type.


\begin{theo}\label{th43}
\source{riemann-roch-polynomials-hyperkahler-fourfolds}
Let $X$ be a  \hk\ fourfold with classes~$\lll,\mm\in H^2(X,\Z)$ such that \mbox{$\int_X \lll^4=0$}.\   Assume $a\coloneqq \frac{1}{2}\int_X \lll^{2}\mm^{2}\in \{2,\dots,8\}$.\
 We are in one of the following cases:
\begin{enumerate}[{\rm (a)}]
 \item\label{enum:th43a} either $a=3$,
\source{riemann-roch-polynomials-hyperkahler-fourfolds}
 $q_X(\lll,\mm)=1$, $c_X=9$,     $P_{RR,X}(T)= 3\binom{\frac{T}2+2}{2}$, and the quadratic form $q_X$ is even; moreover,
\subitem-- either  $(b_2(X),b_3(X),b_4(X))=(7,8, 108)${\rm ;}
 \subitem-- or $(b_2(X),b_3(X),b_4(X))=(6,4, 102)${\rm ;}
\subitem-- or $(b_2(X),b_3(X),b_4(X))=(5,0, 96)$.
\item\label{enum:th43b} or $a=4$, the Chern and Hodge numbers of $X$ are those of the Hilbert square of a K3 surface, and
\source{riemann-roch-polynomials-hyperkahler-fourfolds}
\subitem-- either $q_X(\lll,\mm)=2$, $c_X=3$, the form $q_X$ is even, and $P_{RR,X}(T)=\binom{\frac{T}2+3}{2}${\rm ;}
\subitem-- or $q_X(\lll,\mm)=1$, $c_X=12$, and $P_{RR,X}(T)=\binom{T+3}{2}$.
 \end{enumerate}
\end{theo}

\begin{proof}
We follow the proof of Theorem~\ref{th43a}, which dealt with the case $a=1$:  we may assume  $\gamma\coloneqq \frac{q_X(\mm)}{q_X(\lll,\mm)}\in (-1,1] $, we introduce the polynomial
$$P(k)\coloneqq P_{RR,X}(q_X(k\lll+ \mm))=P_{RR,X}(2kq_X(\lll,\mm)+q_X(\mm)),$$
and we compute
\begin{align*}
 P(k)&=\frac{a}{2}k^2+\Bigl( \frac{a}{2}\gamma+2\sqrt{2aA_X} \Bigr)k+\frac{a}{8}\gamma^2+ \gamma\sqrt{2aA_X}+3\\
 &=:\frac{a}{2}k^2 +  b  k +c.
\end{align*}
Since $P$ takes integral values on integers, $\frac{a}2+b=P(1)-P(0)$ and $c=P(0)$ are integers; when~$a$ is odd, we write $b=\frac12+b'$, with $b'\in\Z$.\ We also check
 \begin{equation}\label{abc}
\source{riemann-roch-polynomials-hyperkahler-fourfolds}
    4A_X-\frac{b^2}{2a} =3-c \in\Z.
\end{equation}
 Finally, when $a$ is not a perfect square, the fact that $\sqrt{2aA_X}$ is rational (Lemma~\ref{lemmax}), implies $A_X\ne\frac{25}{32}$ hence, by Lemma~\ref{lemmguan}\ref{bax},    $\frac56\le A_X\le\frac{131}{144}$.



 By~\eqref{abc}, we  have $4A_X-\frac{b^2}4\in\Z$ and $b\in\Z$.\ By Lemma~\ref{lemmguan}, this case does not happen.


 By~\eqref{abc}, we  have $4A_X- \frac{b^2}6=4A_X-\frac{b'(b'+1)}6-\frac1{24}\in\Z$, hence $4A_X -\frac1{24}\in \Z+\{0,\frac13\}$.\
Since $a$ is not a perfect square, we have (as noted above) $\frac{5}{6}\le A_X\le\frac{131}{144} $, hence
\[
\frac{10}{3} -\frac1{24} \le 4A_X -\frac1{24}\le\frac{131}{36} -\frac1{24}.
\]
The only possibility is  $4A_X-\frac1{24}=3+\frac{1}3$, that is, $A_X=\frac{27}{32}$.\
This implies
$$\frac12+b'=b=\frac{3}{2}\gamma+\frac92,$$
so that $\gamma$ is an even integer.\ As above, this implies $\gamma=q_X(\mm)=0$, $b=\frac92$, and $c = 3$.\ Furthermore,
\[
P_{RR,X}(2kq_X(\lll,\mm))=  \frac{3}{2}k^2 +  \frac92  k +3=3\binom{k+2}{2}.
\]
By Lemma~\ref{star} (applied with $c=2$, $c'=3$, and $q=2q_X(\lll,\mm)$),
we get $q_X(\lll,\mm)=1$, the form~$q_X$ is even, $c_X=9$, and $P_{RR,X}(T)=   3\binom{\frac{T}2+2}{2}$, as in the generalized Kummer case.


From~\eqref{ax}, we obtain $c_4(X)=108$ and $4b_2(X)-b_3(X)=20 $.\
The possible Betti numbers listed in the theorem then follow from \cite{gua}.


 By~\eqref{abc}, we   have $4A_X-\frac{b^2}8\in\Z$ and $b\in\Z$.\
If $\frac{5}{6}\le A_X\le\frac{131}{144} $, we have $b\equiv 2\pmod4$ and $A_X=\frac78$, which contradicts the fact that $\sqrt{2aA_X}$ is   rational.\
Hence, by Lemma~\ref{lemmguan}, we have $A_X=\frac{25}{32}$ and $b^2\equiv 1\pmod8$, so that $b$ is odd.\ Furthermore, $b=2\gamma +5$, hence $\gamma$ is an integer, which can only be $0$ or $1$.

When $\gamma=q_X( \mm)=0$ and $b=5$, we
obtain $c=3$ and
\[
P_{RR,X}(2kq_X(\lll,\mm))=P(k)= 2k^2 + 5  k +3=\binom{2k+3}{2}.
\]
When
$\gamma=1$ and $b=7$, we   get $q_X(\lll,\mm)=q_X(\mm)$ and $c=6$, and
\[
P_{RR,X}((2k+1)q_X(\lll,\mm))=P(k)= 2k^2 +7  k +6=\binom{(2k+1)+3}{2}.
\]
 In both cases, we have
 $$
 P_{RR,X}(q_X(\lll,\mm)T)=\binom{T+3}{2}.
 $$
 By Lemma~\ref{star} (applied with $c=2$, $c'=1$, and $q=q_X(\lll,\mm)$),
we get $q_X(\lll,\mm)\in\{1,2\}$.

Finally, since $A_X=\frac{25}{32}$, Lemma~\ref{lemmguan} implies that  the Chern and Hodge numbers of $X$ are those of the Hilbert square of a K3 surface.


 We   obtain $4A_X-\frac{b^2}{10}\in\Z$ and $b\in  \Z +\frac12$ and, since $a$ is not a perfect square, $\frac{5}{6}\le A_X<1 $.\
We get $A_X=\frac{29}{32} $, which contradicts  $\sqrt{2aA_X}\in\Q$.


 We   obtain $4A_X-\frac{b^2}{12}\in\Z$ and $A_X\in\{\frac56,\frac{15}{16}\}$, which contradicts  $\sqrt{2aA_X}\in\Q$.


 We   obtain $4A_X-\frac{b^2}{14}\in\Z$ and $A_X\in\{\frac{193}{224},\frac{217}{224}\}$, which contradicts  $\sqrt{2aA_X}\in\Q$.


 We   obtain $4A_X-\frac{b^2}{16}\in\Z$ and $A_X=\frac{57}{64}$, which contradicts  $\sqrt{2aA_X}\in\Q$.
\end{proof}

As we saw in Example~\ref{ex:Kumn},  the first item of the case $a=3$ in Theorem~\ref{th43} is realized by \hk\ fourfolds  of generalized Kummer deformation type;  we do not know whether the other two items occur.\
In the case $a=4$, the first item is realized  on \hk\ fourfolds of $\mathrm{K3}^{[2]}$ deformation type with $\mm$ divisible by 2.

Under a condition on some generalized Fujiki constants of $X$, Beckmann and Song prove in~\cite{bs} (see also~\cite{saw2}) that the only possible Betti numbers $b_2(X)$, $b_3(X)$, and $b_4(X)$ for $X$ are as in case~\ref{enum:th43a} of Theorem~\ref{th43}.\
Moreover, by~\cite[Proposition 5.6]{bs}, \cite[Conjecture 1.2]{bs} would imply that the Fujiki constant $c_X$ is either~$3$ or~$9$.\ In particular, the second case in~\ref{enum:th43b} should not occur.



%%%%%%%%%%%%%%%%%%%%%
